% D.91 candidate. Notes only; not registered. Full independent audits pending.
\subsection{The compact kernel for every completed cusp parameter}

\begin{statement}{D.91}{Open}{Candidate: the complete compact-kernel classification on the three cusp packets}
Let
\[
 a=A\kappa(5)+\sum_{q\ne5}a_q\kappa(q),\qquad
 A,a_q\in\mathbb Z_5,\qquad a_q\longrightarrow0,
                                                               \tag{D.91.1}
\]
and consider the actual cusp pair \(s=(a,0)\). The assertion
includes \(a=0\). If all away-five coefficients vanish, then
\(\mathcal S_s=\mathcal K_s=0\), for every \(A\), including
\(A=0\). Thus the trivial Kummer parameter, represented by
the rational tangents \(1\) and \(-1\), is included separately.
Suppose henceforth that the away-five class is nonzero, and put
\[
 \begin{gathered}
 r=\min_{q\ne5}v_5(a_q),\qquad m=v_5(A)\in
                  \mathbb Z_{\ge0}\cup\{\infty\},\\
 \lambda=\sum_{q\ne5}a_q\frac{\log(q^4)}{4\log6},\qquad
 N=v_5(\lambda)\ge r,\qquad e=N-r.
 \end{gathered}                                                \tag{D.91.2}
\]
If \(e=0\), again \(\mathcal S_s=\mathcal K_s=0\).
If \(0<e<\infty\), define
\[
 H=\min\{e,\max(m-r,0)\},\qquad t=\min(m,r),
 \quad H=e\text{ when }m=\infty.
                                                               \tag{D.91.3}
\]
Then, as compact topological modules,
\[
 \begin{aligned}
 \mathcal S_s&\simeq
 \left(\prod_{n\ge0}\mathbb Z/5^e\right)
 \times\prod_{j=1}^{\min(t,e-H-1)}
                      (\mathbb Z/5^{H+j})^2,\\
 \mathcal K_s&\simeq\mathcal S_s\times\mathbb Z_5^{\,2H}.
 \end{aligned}                                                \tag{D.91.4}
\]
A finite product with upper limit less than one is zero.
If \(N=\infty\) and \(m<\infty\), put
\(h=\max(m-r,0)\) and \(t=\min(m,r)\). Then
\[
 \begin{aligned}
 \mathcal S_s&\simeq
 \left(\prod_{n\ge0}\mathbb Z_5\right)
          \times\prod_{j=1}^{t}(\mathbb Z/5^{h+j})^2,\\
 \mathcal K_s&\simeq\mathcal S_s\times\mathbb Z_5^{\,2h}.
 \end{aligned}                                                \tag{D.91.5}
\]
If \(N=m=\infty\), both \(\mathcal S_s\) and
\(\mathcal K_s\) are isomorphic to \(\prod_{n\ge0}\mathbb Z_5\).

The subgroup \(\mathcal S_s\subset\mathcal K_s\) and
the map \(L_5\to\mathcal T_s\) are the actual natural
maps of D.63. The displayed splittings and identifications
with products are continuous and noncanonical. In (D.91.5)
one may absorb the finite free factor into the countable
free product as an abstract compact-module isomorphism;
this does not identify that factor with the natural kernel
\(\mathcal S_s\).

The same classification applies to all three peripheral
markings: the actual tripod automorphisms transport
\((a,0)\) to \((0,a)\) and to \((-a,-a)\), together
with the tangential frame. They induce isomorphisms of the
entire global/local cohomology diagrams. In particular no
zero-parameter packet at one of the other two cusps is
omitted or identified with another section packet.

Every parameter above has an actual cusp section. For finite
indices the rational tangent
\[
                5^{5^m}(1+5^{e+1})^{5^r}
                                                               \tag{D.91.6}
\]
realizes every \(r,m\ge0\) and \(e\ge1\); omit the first
factor for \(m=\infty\). For example, \(5^{25}626^5\)
has \((r,m,N,e)=(1,2,4,3)\), and hence
\[
 \mathcal S_s\simeq
 \left(\prod_{n\ge0}\mathbb Z/125\right)\times(\mathbb Z/25)^2,
 \qquad \mathcal K_s\simeq\mathcal S_s\times\mathbb Z_5^2.
                                                               \tag{D.91.7}
\]
Their finitely ramified metabelian fibres have translation
group \(\prod_{\mathbb N}\mathbb Z_5\), by D.74.
The theorem computes the compact second-cohomology kernel
on the cusp families, not every abelian pair, every
infinitely ramified first-cohomology rank, or a first
higher-derived obstruction.
\end{statement}

\begin{proof}[Candidate proof, pending complete independent audits]
Write \(\Lambda=\mathbb Z_5[[u,v]]\), \(t=1+u\),
\(w=t(1+v)-1\), \(P=uvw\), and \(q_j=t^{5^j}-1\),
including \(q_0=u\). Every quotient below is compact,
with its quotient topology. All actions are the actual
operator actions of D.35, D.39 and D.87.

\paragraph{Degeneracies and the ranges already computed.}
If the away-five class is zero, the actual global and local
images are the same graph of \(A\kappa(5)\) over \(I_0\).
D.83 and D.87 give \(\mathcal S_s=0\), and D.77 gives
\(\mathcal K_s=0\). This includes \(a=0\); it is not
obtained by subtracting two infinite valuations.
For nonzero away class the actual images are
\[
 I_s=5^r\mathbb Z_5\rtimes I_0,\qquad
 I_5=5^N\mathbb Z_5\rtimes I_0,
                                                               \tag{D.91.8}
\]
with the same graph action
\(\rho_A(g)=t^{A\kappa_5(g)}\rho_0(g)\).
Here \(5^\infty\mathbb Z_5=0\). D.83's uniformly
convergent integral gauge belongs to the actual global
translation lattice, so (D.91.8) holds also for infinite
\(c_0\) support. If \(N=r\), the images coincide.

For \(A=0\), D.88 proves the asserted answer, including
\(N=\infty\). For finite \(m\le r\), D.90 proves it.
If \(r\le N\le m<\infty\), the graph multiplier is
already one on \(R=\Lambda/(q_N)\) and on \(J=q_rR\).
The entire natural invariant map in D.87 is therefore
literally the D.88 map, not just a map between isomorphic
abstract modules. D.88 computes \(\mathcal S_s\), and
D.77 adds its free factor \(\mathbb Z_5^{2(N-r)}\).
It remains to prove the genuinely new range
\[
                 0\le r<m<N\le\infty.
                                                               \tag{D.91.9}
\]
Put \(h=m-r>0\). For finite \(N\), write \(b=N-m>0\),
so \(e=h+b\).

\paragraph{The actual scalar module and the map that must be kept.}
The two actual operators of D.89/D.90 give
\[
       (\rho_A(g)-1:g\in\ker\chi)=(P,q_m).
                                                               \tag{D.91.10}
\]
Let \(R=\Lambda/(q_N)\) for finite \(N\), and \(R=\Lambda\)
otherwise; put \(J=q_rR\). D.87 identifies \(\mathcal S_s\)
with the kernel of the natural compact map
\[
 M:=\bigl(J/(PJ+q_mJ)\bigr)(-1)_\Gamma
       \longrightarrow
 U:=\bigl(R/(PR+q_mR)\bigr)(-1)_\Gamma,
 \qquad \Gamma=\mathbb Z_5^\times.
                                                               \tag{D.91.11}
\]
The map is induced by \(J\hookrightarrow R\).
In particular it is not discarded when \(r<m\).
The target is the ordinary dual of the actual local
Cartier-invariant group, so D.77 identifies it with
\(L_5\); it is torsion-free.

Retain \(q_r\) as the actual generator of \(J\).
For finite \(N\), set \(F=q_m/q_r\). The exact norm
identity modulo \(q_m\) gives
\[
 q_N/q_r\equiv5^b F\pmod{q_m},\qquad
 B=\mathbb Z_5[[u,v]]/(q_m,5^bF),\qquad T=B/(P).
                                                               \tag{D.91.12}
\]
For \(N=\infty\), omit the relation \(5^bF\) and put
\(B=\mathbb Z_5[\mathbb Z/5^m][[v]]\), again \(T=B/(P)\).
After the Cartier untwist, the action on the scalar module
\(T\) is
\[
 \chi(g)^{-1}S_{\chi(g)}(t^{5^r})\rho_0(g).
                                                               \tag{D.91.13}
\]
The graph multiplier is one modulo \(q_m\). Thus
\(M=T_\Gamma\) with precisely (D.91.13). The target
has scalar coefficient \(\Lambda/(q_m,P)\), action
\(\chi^{-1}\rho_0\), and the map in (D.91.11) is
\emph{multiplication by \(q_r\)}.

\paragraph{The infinite layer injects before any boundary calculation.}
Let \(C=\mathbb Z/5^e\) for finite \(N\), and
\(C=\mathbb Z_5\) for \(N=\infty\). Since \(F(0)=5^h\),
we have \(B/uB=C[[v]]\). Both \(v\) and \(w\) are
nonzero divisors in \(B\): multiplication by \(v\) is
injective, and the continuous change of variables
\(v\mapsto u+(1+u)v\) in the complete \((5,u,v)\)-adic
ring carries it to multiplication by \(w\). Consequently
\[
 0\longrightarrow E\longrightarrow T\longrightarrow Q\longrightarrow0,
 \qquad E=vwB/(uvwB)\simeq C[[v]],\quad Q=B/(vw).
                                                               \tag{D.91.14}
\]
The retained generator \(q_rvw\) gives on \(E\) the
action \(\chi S_\chi(1+v)^3\) followed by substitution.
D.88's exact weighted-shell calculation gives
\(E_\Gamma\simeq\prod_{n\ge0}C\).

The two successive tame characters of \(Q/uQ\) are
\(\chi^2\) and \(\chi^3\). Therefore every invariant
element of \(Q\) lies in \(uQ\). Moreover
\(uT\to uQ\) is an equivariant isomorphism: modulo
\(u\), the element \(vw\) becomes the nonzero divisor
\(v^2\) in \(C[[v]]\), so \((u)\cap(vw)=(uvw)\).
Every invariant lifts invariantly. The snake lemma for
the actual lift of \(2\in\Gamma\) yields the exact sequence
\[
       0\longrightarrow E_\Gamma\longrightarrow M
                 \longrightarrow Q_\Gamma\longrightarrow0.
                                                               \tag{D.91.15}
\]
The same proof applies to coefficient reductions modulo
\(5^H\) in the infinite case. This is an actual injection,
not a deduction from a comparison of ranks.

\paragraph{Axis intersections and the literal boundary modules.}
Write \(\mathcal N_m=q_m/u\). In the finite case
\[
                 \operatorname{Ann}_B(u)
                   =(\mathbb Z/5^b)[[v]]\mathcal N_m.
                                                               \tag{D.91.16}
\]
Indeed, if \(uf=F(q_rg+5^b k)\) in the original
torsion-free power-series ring, evaluation at \(u=0\)
forces \(k(0,v)=0\). Cancel \(u\) to express \(f\),
modulo the relations, as \(\mathcal N_m g\). Since
\(u\mathcal N_m=0\), only \(g(0,v)\) remains. Its
coefficient has exact modulus \(5^b\): cancel \(F\)
and reduce modulo \(q_r\), where the norm
\(q_r/u\) has every cyclic coefficient equal to one.
For \(N=\infty\), the same literal cyclic coefficient
argument gives \(\operatorname{Ann}_B(u)=\mathbb Z_5[[v]]\mathcal N_m\).

The two-axis map for \(Q\) consequently has kernel
\(Cv\mathcal N_m\), with \(C\) here replaced by
\(\mathbb Z/5^b\) in the finite case, and cokernel
\(B/(u,v)\). The first has exact tame weight \(\chi^3\),
the second weight \(\chi^2\). To see the kernel directly,
write \(vf=wg\), reduce modulo \(v\), use (D.91.16),
and note \(w\mathcal N_m=v\mathcal N_m\) and
\(v^2\mathcal N_m=0\) in \(Q\). Both kernel and
cokernel have zero invariants and coinvariants under the
order-four tame subgroup. They therefore contribute no
\(\Gamma\)-homology.

Let \(R_j=\mathbb Z_5[\mathbb Z/5^j]\) and
\(J_{r,N}=\ker(R_N\to R_r)\). On the first axis,
(D.91.13) is
\(S_\chi(t)S_\chi(t^{5^r})\sigma\), with
\(\sigma(t)=t^\chi\). Multiplication by \(u q_r\)
transports it exactly to the ordinary permutation action on
\(uJ_{r,N}/q_m uJ_{r,N}\). Multiplication by \(u\)
is injective on \(J_{r,N}\): its only annihilator in
\(R_N\) is the full norm, which has nonzero image in
the torsion-free ring \(R_r\). The comparison
\[
       uJ_{r,N}/q_m uJ_{r,N}
                      \longrightarrow J_{r,N}/q_mJ_{r,N}
                                                               \tag{D.91.17}
\]
has kernel \(q_mJ_{r,N}/u q_mJ_{r,N}\), of tame
weight \(\chi^2\), and cokernel \(J_{r,N}/uJ_{r,N}\),
of tame weight \(\chi\). They vanish in all
\(\Gamma\)-homology. The diagonal axis has the same
comparison after multiplication by \(t^{-1/2}\), from
the exact reflection identity of D.39. Thus
\[
 Q_\Gamma\simeq
          \left((J_{r,N}/q_mJ_{r,N})_\Gamma\right)^2.
                                                               \tag{D.91.18}
\]
At \(N=\infty\), use \(J_{r,\infty}=q_r\mathbb Z_5[[t-1]]\);
the same argument and tame weights apply.

\paragraph{The middle boundary is a nonsplit cyclic extension.}
For finite \(N\), let \(D=\mathbb Z/5^b\) act on
\(R_N\) by translation by \(5^m\), and let
\(\mathcal Q=D\rtimes\Gamma\). Its action on \(R_r\)
has trivial \(D\)-part. Compact group homology of
\(0\to J_{r,N}\to R_N\to R_r\to0\) gives
\[
 H_1(\mathcal Q,R_N)\longrightarrow H_1(\mathcal Q,R_r)
 \longrightarrow (J_{r,N})_{\mathcal Q}
 \longrightarrow (R_N)_{\mathcal Q}\longrightarrow(R_r)_{\mathcal Q}
 \longrightarrow0.                                           \tag{D.91.19}
\]
These are permutation modules on finite sets. For complete
control of continuity and the map, dualize to the discrete
modules of functions valued in \(\mathbb Q_5/\mathbb Z_5\).
Continuous Shapiro at each open point stabilizer identifies
first cohomology with its continuous characters. Dualizing
back gives the pro-five abelianization of that stabilizer.
The finite-index cochain proof is the usual Shapiro map
\cite[II, Proposition 1.11 and Example 1.27(a)]{PBC_MilneCFT};
it applies to continuous cochains and discrete coefficients.
Pullback of functions along an equivariant point-set map
becomes restriction of characters. Its compact dual is
therefore the stabilizer inclusion, not a transfer. This
also constructs the exact sequence (D.91.19) directly by
dualizing the discrete cohomology sequence.

There are \(m+1\) upper orbits, the valuation shells
\(0,\ldots,m-1\) and the ball divisible by \(5^m\),
and \(r+1\) lower orbits. The kernel on free orbit
coinvariants is \(\mathbb Z_5^h\). For a lower nonzero
orbit represented by \(5^j\), put \(k=r-j\), so
\(1\le k\le r\). Its stabilizer is
\(D\rtimes(1+5^k\mathbb Z_5)\), whose pro-five
abelianization is
\[
             \mathbb Z/5^{\min(b,k)}\oplus\mathbb Z_5.
                                                               \tag{D.91.20}
\]
The sole upper orbit over it has stabilizer the graph
\(d=(1-\chi)/5^{k+h}\), with
\(\chi\in1+5^{k+h}\mathbb Z_5\). Choose its generator
\(\chi=(1+5^k)^{5^h}\). The actual stabilizer inclusion
maps it to \((d,5^h)\), and \(d\) is a five-adic unit.
Its cokernel in (D.91.20) is therefore a \emph{single}
cyclic group of order \(5^{h+\min(b,k)}\).
For the lower zero orbit, tame conjugation kills the
translation contribution to the abelianization, and the
upper zero-orbit stabilizer \(\Gamma\) surjects on the
remaining first homology. There is no contribution there.
It follows that
\[
 (J_{r,N}/q_mJ_{r,N})_\Gamma
       \simeq\mathbb Z_5^h\oplus
                       \bigoplus_{k=1}^r\mathbb Z/5^{h+\min(b,k)}.
                                                               \tag{D.91.21}
\]
The finite free quotient splits continuously. In particular
\((r,m,N)=(1,2,3)\) gives \(\mathbb Z_5\oplus\mathbb Z/25\),
not \(\mathbb Z_5\oplus(\mathbb Z/5)^2\).

\paragraph{The actual boundary map has exactly that torsion kernel.}
On the target first axis the untwisted action is
\(S_\chi(t)\sigma\). Multiplication by \(u\) transports
it to ordinary permutation on \(uR_m\), with a kernel
of tame weight \(\chi\). The map \(uR_m\to R_m\)
is injective on coinvariants: the fixed element \(\delta_0\)
makes \(H_1(\Gamma,R_m)\to H_1(\Gamma,\mathbb Z_5)\)
surjective. Its image is the kernel of augmentation in
the free orbit module \((R_m)_\Gamma\).

The source and target transports commute with the actual
map \(f\mapsto q_rf\): both composites are multiplication
by \(u q_r\). The diagonal half-gauges commute with
\(q_r\) as well. Thus the induced boundary map is
exactly the literal reduction and inclusion
\[
 (J_{r,N}/q_mJ_{r,N})_\Gamma\longrightarrow(R_m)_\Gamma.
                                                               \tag{D.91.22}
\]
In (D.91.19), it kills the finite homology cokernel and
identifies the free orbit-kernel quotient with
\(\ker((R_m)_\Gamma\to(R_r)_\Gamma)\). Its kernel
is precisely the finite torsion in (D.91.21), with no
kernel on the free quotient. The target axis-intersection
kernel and cokernel have tame weights \(\chi^2\) and
\(\chi\), so the same assertion holds for both actual
boundary axes. No comparison by abstract coranks is used.

\paragraph{Finite depth: the entire extension is killed before it is split.}
Let \(F_b\) denote two copies of the finite torsion in
(D.91.21). Equations (D.91.15) and (D.91.22) show that
the kernel in (D.91.11) is contained in the preimage
of \(F_b\). That preimage is exactly the torsion subgroup
of \(M\), because \(E_\Gamma\) has exponent \(5^e\)
and the remaining quotient is finite free. Conversely all
this torsion maps to zero in the torsion-free actual target
\(U\). Therefore
\[
  \mathcal S_s=M_{\rm tor},\qquad
 0\longrightarrow\prod_{n\ge0}\mathbb Z/5^e
   \longrightarrow\mathcal S_s\longrightarrow F_b\longrightarrow0.
                                                               \tag{D.91.23}
\]
The finite-index actual images in (D.91.8) have index
\(5^e\). Restriction followed by corestriction on their
actual continuous first cohomology, with its actual
coefficient action, is multiplication by \(5^e\).
D.88 gives the explicit finite-coset proof, including
nonnormal subgroups; see also \cite[II, Proposition 1.30]{PBC_MilneCFT}.
Consequently the \emph{whole} group \(\mathcal S_s\)
in (D.91.23) is killed by \(5^e\).

For a cyclic quotient generator of order \(5^c\), choose
a lift \(x\). Then \(5^cx\) lies in the product layer
and is annihilated by \(5^{e-c}\). In that product,
\(\operatorname{Ann}(5^{e-c})=5^c E_\Gamma\). Subtract
an element of the layer to make the lift have order
\(5^c\). Lifting the finitely many cyclic generators
this way splits (D.91.23) continuously. The factors with
\(h+\min(b,k)=e\) are absorbed by the countable product
of \(\mathbb Z/5^e\); the smaller factors are retained.
This proves the strict-middle instance of (D.91.4).

\paragraph{The infinite boundary and its actual reduction maps.}
Keep \(m,r\) fixed and let the finite integer \(b\)
increase. The scalar boundary modules are
\(R_m/(5^bF)\) with the fixed retained boundary action
\(\varphi_r=S_\chi(t)S_\chi(t^{5^r})\sigma\). Their
inverse limit is \(R_m\): their defining ideals tend
uniformly to zero. Compactness also gives
\[
  (R_m,\varphi_r)_\Gamma
       \simeq\varprojlim_b(R_m/(5^bF),\varphi_r)_\Gamma.
                                                               \tag{D.91.24}
\]
Explicitly, compact intersections give surjectivity onto
compatible quotient cosets, and
\(\bigcap_b((2-1)R_m+5^bFR_m)=(2-1)R_m\), where
\(2-1\) denotes the actual operator difference. Its
image is compact and closed; the ideals tend to zero.

The stabilizer maps in (D.91.20) reduce the translation
coordinate and leave the free orbit-kernel map unchanged.
For \(b\ge r\), all finite cyclic factors have stabilized,
and the transition maps are isomorphisms: they are onto,
are the identity on the free orbit quotient, and have
equal finite torsion orders. Hence the infinite boundary is
\[
 (J_{r,\infty}/q_mJ_{r,\infty})_\Gamma
       \simeq\mathbb Z_5^h\oplus
                              \bigoplus_{k=1}^r\mathbb Z/5^{h+k}.
                                                               \tag{D.91.25}
\]
Its natural map to \((R_m)_\Gamma\) still kills exactly
the finite torsion and embeds the free orbit part, by
the same actual maps, or by their stabilized limit.

\paragraph{Integral torsion lifts at infinity, and the actual kernel.}
Now \(E_\Gamma=\prod\mathbb Z_5\). Put
\(F_\infty=\bigoplus_{k=1}^r(\mathbb Z/5^{h+k})^2\).
Split the finite free orbit quotient of \(M\), and let
\(M_0\) be the preimage of \(F_\infty\). Then
\[
 0\longrightarrow E_\Gamma\longrightarrow M_0
              \longrightarrow F_\infty\longrightarrow0,
 \qquad M\simeq M_0\times\mathbb Z_5^{2h}.
                                                               \tag{D.91.26}
\]
For every \(H\ge1\), reduction of the actual scalar
module modulo \(5^H\) commutes with coinvariants. The
zero-snake injection proved after (D.91.15) remains valid
over \(\mathbb Z/5^H\). It gives the exact purity identity
\[
             E_\Gamma\cap5^H M=5^H E_\Gamma.
                                                               \tag{D.91.27}
\]
For \(H\ge m\), the quotient \(F_\infty\) is killed
by \(5^H\), so \(5^HM_0\subset E_\Gamma\). Equation
(D.91.27) therefore gives \(5^HM_0=5^HE_\Gamma\).
The finite-coefficient extension
\(M_0/5^HM_0\) is killed by \(5^H\), has the product
layer \(E_\Gamma/5^HE_\Gamma\), and quotient
\(F_\infty\). The same finite-order adjustment as above
provides a lift of each quotient generator of its exact
order modulo \(5^HM_0\). Lift it integrally. If its
order is \(5^c\), its \(5^c\)-multiple lies in
\(5^HM_0=5^HE_\Gamma\subset5^cE_\Gamma\); one more
integral adjustment gives an actual order-\(5^c\) lift.
Thus (D.91.26) splits continuously, retaining its finite
torsion rather than inferring a splitting from abstract ranks.

The actual map in (D.91.11) kills the whole product layer:
its actual generator maps as
\[
                      q_rvw=P\,(q_r/u)=0
                          \quad\text{in }\Lambda/(P,q_m).
                                                               \tag{D.91.28}
\]
It also kills the finite torsion, because its target is
torsion-free. Conversely (D.91.22)/(D.91.25) embeds the
free orbit quotient, so nothing outside \(M_0\) can
lie in its kernel. Therefore
\(\mathcal S_s=M_0\simeq E_\Gamma\times F_\infty\).
This proves the strict-middle case of (D.91.5) by actual
maps, including the simultaneous absence of any local
unit translation.

\paragraph{The remaining boundary factor, all cusp markings, and examples.}
D.77 identifies the zero-parameter boundary map with the
identity on its countable product. The other two boundary
depths are globally \(\min(m,r)\) and locally
\(\min(m,N)\). For finite \(N\) their difference is
exactly \(H\) in (D.91.3); for infinite \(N\) and
finite \(m\) it is \(h\). This gives the natural compact
quotient \(\mathbb Z_5^{2H}\), or \(\mathbb Z_5^{2h}\),
and the continuous noncanonical split sequence of D.63.
When \(m=N=\infty\), D.88 gives the stated product
identifications with their precise natural-map interpretation.

The actual automorphism \(t\mapsto1-t\) exchanges the
first two cusp packets. The actual automorphism
\(t\mapsto1/t\) transforms the Kummer pair as
\((a,b)\mapsto(-a,b-a)\), because \(\kappa(-1)=0\)
in pro-five Kummer cohomology. It therefore sends
\((a,0)\) to \((-a,-a)\). Transport the tangential
frame along the automorphism; changing back to the standing
frame is an inner gauge. These are actual isomorphisms of
the twisted Alexander modules and of every finite global
and local cohomology map. Replacing \(a\) by \(-a\)
does not change \(m,r,N\). This proves the asserted
three-packet scope, without identifying their sections.

For (D.91.6), \(1+5^{e+1}\equiv2\pmod4\), so
its exponent at two is exactly one and its Kummer class
is globally primitive. Its logarithmic unit coefficient
has valuation \(e\). Raising it to \(5^r\) gives
the required \(r,N=r+e\), and the independent five-power
factor supplies \(m\). For \(N=\infty\), use
\(5^r(\lambda_3\kappa(2)-\lambda_2\kappa(3))\),
where \(\lambda_2\equiv2\) and \(\lambda_3\equiv4\pmod5\),
and add \(5^m\kappa(5)\) if desired. All are actual
completed cusp parameters, with the full cusp construction
supplying section existence. D.74 supplies the asserted
finitely ramified fibre statement.

The literal scalar matrices, norm relations, actual
source-to-target maps and diagonal gauge are checked in
\path{notes/de_sha_cusp_complete_certificate.py}.
The independent stabilizer calculation is retained in
\path{notes/pbc_middle_boundary_derivation.md} and its
exact certificate. Those finite computations support the
formulas; the all-depth compact and arithmetic arguments
are the proof above.
\end{proof}
